\originalchapter{小振动与正常模}
\originalsection{稳定平衡附近的二次型}
设固定完整约束系统在$q_0$处静止平衡, 势满足$\nabla U(q_0)=0$. 取小偏离$\eta=q-q_0$, 忽略三阶及以上小量, 写$T=\dot\eta^{\mathsf T}M\dot\eta/2$, $U=U(q_0)+\eta^{\mathsf T}K\eta/2$. $M$是平衡点处的质量矩阵, 对独立坐标正定; $K$是势能Hessian, 是实对称矩阵. 线性运动方程为$M\ddot\eta+K\eta=0$.

令$\eta=a e^{i\omega t}$, 得$Ka=\omega^2Ma$. 这里不是直接对$K$作普通对角化, 因为坐标的惯性度量是$M$. 取对称正定平方根$M^{1/2}$, 令$B=M^{-1/2}KM^{-1/2}$, 则$B$实对称. 谱定理给正交基$u_\alpha$和实特征值$\omega_\alpha^2$. 定义$a_\alpha=M^{-1/2}u_\alpha$, 有$a_\alpha^{\mathsf T}Ma_\beta=\delta_{\alpha\beta}$. 写$\eta=\sum_\alpha a_\alpha Q_\alpha$, 便得到
\[
L_2=\frac12\sum_\alpha(\dot Q_\alpha^2-\omega_\alpha^2Q_\alpha^2),\qquad
\ddot Q_\alpha+\omega_\alpha^2Q_\alpha=0.
\]
若$K$正定则频率全实且非零. 若有零特征值, 二次近似中该方向没有回复力. 它可能来自自由移动或转动, 也可能是二阶项偶然消失而高阶项仍有回复力, 例如$U=x^4$的原点. 此时线性模型不能给有限非零振动频率, 需检查高阶项. 若有负特征值, 微扰指数增长, 平衡不稳定. 重根不妨碍选取$M$正交的模, 但重根子空间内的基不唯一.

\originalsection{双摆的线性化}
\begin{example}
两个无质量杆长均为$\ell$, 上端固定, 质量$m_1,m_2$分别在第一杆末端和第二杆末端. 两杆相对竖直向下的绝对角为$\theta_1,\theta_2$. 求$m_1=3m$, $m_2=m$时的小振动.
\end{example}
\begin{solution}
第一质量位置为$(\ell\sin\theta_1,-\ell\cos\theta_1)$, 第二质量再加$(\ell\sin\theta_2,-\ell\cos\theta_2)$. 速度平方展开给
\[
T=\tfrac12(m_1+m_2)\ell^2\dot\theta_1^2+\tfrac12m_2\ell^2\dot\theta_2^2
+m_2\ell^2\dot\theta_1\dot\theta_2\cos(\theta_1-\theta_2).
\]
势为$U=-(m_1+m_2)g\ell\cos\theta_1-m_2g\ell\cos\theta_2$. 二次近似的矩阵是$M=m\ell^2\left(\begin{smallmatrix}4&1\\1&1\end{smallmatrix}\right)$, $K=mg\ell\left(\begin{smallmatrix}4&0\\0&1\end{smallmatrix}\right)$. 令$\lambda=\omega^2\ell/g$, 特征行列式为$(4-4\lambda)(1-\lambda)-\lambda^2=3\lambda^2-8\lambda+4$. 两根是$2/3$和$2$, 均为正. 相应方向为$(1,2)$与$(1,-2)$: 低频模同向, 高频模反向. 若取$M$归一化, 分别除以$\sqrt{12m\ell^2}$和$\sqrt{4m\ell^2}$. 任意小初值都可投影为$Q_\alpha(0)=a_\alpha^{\mathsf T}M\theta(0)$和$\dot Q_\alpha(0)=a_\alpha^{\mathsf T}M\dot\theta(0)$, 再逐模求解.
\end{solution}

\originalsection{阻尼与受迫振动}
线性阻尼力$-b\dot x$把机械能转为其他形式. 对$m\ddot x+b\dot x+kx=f(t)$, 定义$\omega_0^2=k/m$, $\gamma=b/(2m)$. 齐次根为$-\gamma\pm\sqrt{\gamma^2-\omega_0^2}$. 欠阻尼时振幅乘$e^{-\gamma t}$, 振荡频率为$\sqrt{\omega_0^2-\gamma^2}$; 临界阻尼有$(A+Bt)e^{-\gamma t}$; 过阻尼为两个不同实指数. 若$b,k>0$, 各根实部均负, 初始扰动衰减.

若$f=F\cos\omega t$, 用$\Re(Xe^{i\omega t})$求稳态, 则$X=F/(k-m\omega^2+ib\omega)$. 因而幅值为$F/\sqrt{(k-m\omega^2)^2+b^2\omega^2}$, 相位滞后为$\operatorname{atan2}(b\omega,k-m\omega^2)$. 阻尼非零时“幅值最大”频率为$\sqrt{\omega_0^2-2\gamma^2}$（当根存在）, 与无阻尼固有频率不完全相同. 若$b=0$且$\omega=\omega_0$, 有特解$Ft\sin\omega_0t/(2m\omega_0)$, 振幅线性增长; 该增长会最终破坏原先的小振动假设.

对多坐标且阻尼矩阵不与$M,K$共同对角化的系统, 正常坐标未必把阻尼也分离. 不能仅把单振子的阻尼根逐一贴到保守正常模上. 本册在线性保守模和一维阻尼振子的范围内给出完整推导.

本章来源: 8.223第10--13讲及作业3, 8.09第3章. 双摆是根据系统定义自编的中文例题, 不复制被暂缓打包的商业教材题面.
